\subsection{Operasi Elementer}

Kita memulai bagian ini dengan sebuah contoh. Ingat kembali dari Contoh \ref{exa:graphicalsoln} bahwa solusi sistem yang diberikan adalah $\left(x, y \right) = \left( -1, 4 \right)$.

\begin{example}{Memverifikasi bahwa Pasangan Terurut Merupakan Solusi}\label{exa:verifyingorderedpairsolution}
Verifikasikan secara aljabar bahwa $\left(x, y \right) = \left( -1, 4 \right)$ merupakan solusi sistem persamaan berikut.
\[
\begin{array}{c}
x+y=3 \\
y-x=5  
\label{eq1.2.1}
\end{array}
\]
\end{example}

\begin{solution} 
Dengan menggambarkan kedua persamaan ini dan menentukan titik
perpotongannya, sebelumnya kita menemukan bahwa $\left(x, y \right) = \left(
-1, 4 \right)$ merupakan solusi unik.

Kita dapat memverifikasinya secara aljabar dengan menyubstitusikan nilai-nilai ini
ke dalam persamaan semula dan memastikan bahwa persamaan tersebut berlaku.
Pertama, kita menyubstitusikan nilainya ke dalam persamaan pertama dan memeriksa bahwa hasilnya sama dengan $3$.
\[
x + y = (-1)+(4) = 3
\]
Hasil ini sama dengan $3$ seperti yang diperlukan, sehingga kita melihat bahwa $\left( -1,4 \right)$ merupakan solusi persamaan pertama.
Menyubstitusikan nilai-nilai tersebut ke dalam persamaan kedua menghasilkan
\[
y -x = (4) - (-1) = 4 + 1  = 5
\]
yang benar.
Untuk
$\left( x,y\right) =\left( -1,4\right) $, setiap persamaan benar; oleh karena itu, pasangan tersebut merupakan solusi sistem.
\end{solution}

Sekarang muncul pertanyaan menarik berikut:\ Jika Anda tidak diberi
bilangan-bilangan ini untuk diverifikasi, bagaimana Anda dapat menentukan solusinya secara aljabar? Aljabar linear memberi kita alat yang diperlukan untuk menjawab
pertanyaan tersebut. Operasi-operasi dasar berikut merupakan alat penting yang akan kita gunakan.

\begin{definition}{Operasi Elementer}\label{def:elementaryoperations}
\textbf{Operasi elementer} \index{operasi elementer} adalah
operasi-operasi berikut.

\begin{enumerate}
\item Tukarkan urutan pencantuman persamaan.

\item Kalikan sembarang persamaan dengan bilangan tak nol.

\item Ganti sembarang persamaan dengan persamaan itu sendiri ditambah suatu kelipatan dari persamaan
lain.

\end{enumerate}
\end{definition}

Penting untuk dicatat bahwa tidak satu pun dari operasi ini akan
mengubah himpunan solusi sistem persamaan. Bahkan, operasi elementer merupakan {\em alat utama \em}yang kita gunakan dalam aljabar linear untuk menemukan solusi sistem persamaan.

Perhatikan contoh berikut.

\begin{example}{Dampak Operasi Elementer}\label{exa:effectsofelementaryoperation}
\label{example:effectsofelementaryoperation}
Tunjukkan bahwa sistem
\begin{equation*}
\begin{array}{cccccc}
x&+&y&=&7 \\
2x&-&y&=&8
\end{array}
\end{equation*}
memiliki solusi yang sama dengan sistem
\begin{equation*}
\begin{array}{cccccc}
x&+&y&=&7 \\
&&-3y&=&-6
\end{array}
\end{equation*}
\end{example}

\begin{solution}
Perhatikan bahwa sistem kedua diperoleh dengan mengambil persamaan kedua dari sistem pertama
dan menambahkan -2 kali persamaan pertama, sebagai berikut:
\begin{equation*}
2x-y + (-2)(x+y) = 8 + (-2)(7)
\end{equation*}
Dengan menyederhanakannya, kita memperoleh
\begin{equation*}
-3y=-6
\end{equation*}
yang merupakan persamaan kedua dalam sistem kedua.
Sekarang, dari sini kita dapat mencari $y$ dan melihat bahwa $y=2$. Selanjutnya, kita menyubstitusikan nilai ini ke dalam persamaan pertama sebagai berikut \:
\begin{equation*}
x+y=x+2=7
\end{equation*}
Jadi $x=5$, sehingga $\left( x,y\right) = \left(5,2 \right)$ merupakan solusi sistem kedua.
Kita ingin memeriksa apakah $\left(5,2 \right)$ juga merupakan solusi sistem pertama. Kita memeriksanya dengan menyubstitusikan $\left(x, y \right) = \left(5,2 \right)$
ke dalam sistem dan memastikan bahwa persamaan-persamaannya benar.
\begin{equation*}
\begin{array}{c}
x+y = \left(5 \right)+ \left( 2 \right) = 7 \\
2x-y= 2 \left(5 \right) - \left( 2 \right) = 8
\end{array}
\end{equation*}
Jadi, $\left(5,2 \right)$ juga merupakan solusi sistem pertama.
\end{solution}

Contoh ini memperlihatkan bahwa operasi elementer yang diterapkan pada sistem dua persamaan dalam dua variabel
tidak memengaruhi himpunan solusi. Namun, suatu sistem linear dapat melibatkan banyak persamaan dan banyak variabel, dan tidak ada
alasan untuk membatasi kajian kita pada sistem kecil. Untuk sistem berukuran apa pun dengan berapa pun banyak variabel,
himpunan solusi tetap merupakan kumpulan solusi persamaan-persamaan tersebut. Dalam setiap
kasus, operasi-operasi pada Definisi \ref{def:elementaryoperations} di atas tidak
mengubah himpunan solusi sistem persamaan linear.

Dalam teorema berikut, kita menggunakan notasi $E_i$ untuk merepresentasikan suatu persamaan, sedangkan $b_i$ menyatakan
konstanta.

\begin{theorem}{Operasi Elementer dan Solusi}\label{thm:elementaryoperationsandsolns}
Misalkan Anda memiliki sistem dua persamaan linear
\begin{equation}
 \begin{array}{c}
  E_{1}=b_{1}\\
  E_{2}=b_{2}
\end{array} \label{system}
\end{equation}
Maka sistem-sistem berikut memiliki himpunan solusi yang sama dengan \ref{system}:
\begin{enumerate}
\item   \begin{equation}
	\begin{array}{c}
	E_{2}=b_{2}\\
	E_{1}=b_{1}
	\end{array}
	\label{thm1.9.1}
	\end{equation}
\item  \begin{equation}
	\begin{array}{c}
	E_{1}=b_{1} \\
	kE_{2}=kb_{2}\\        
	\end{array}
	\label{thm1.9.2}
	\end{equation}
  untuk sembarang skalar $k$, asalkan $k\neq0$.
\item \begin{equation}
      \begin{array}{c}
       E_{1}=b_{1} \\
       E_{2}+kE_{1}=b_{2}+kb_{1}
       \end{array}  
	\label{thm1.9.3}
	\end{equation}
	untuk sembarang skalar $k$ (termasuk $k=0$).

\end{enumerate}
\end{theorem}

Sebelum melanjutkan ke pembuktian Teorema \ref{thm:elementaryoperationsandsolns},
marilah kita membahas teorema ini dalam konteks Contoh \ref{exa:effectsofelementaryoperation}. Dalam hal ini,
\begin{equation*}
\begin{array}{cc}
E_{1} = x+y, & b_{1} = 7 \\
E_{2} = 2x-y, & b_{2} = 8 
\end{array}
\end{equation*}
Ingat kembali operasi elementer yang kita gunakan untuk mengubah sistem dalam solusi contoh tersebut.
Pertama, kita menambahkan $\left( -2 \right)$ kali persamaan pertama ke persamaan kedua.
Menurut Teorema \ref{thm:elementaryoperationsandsolns}, tindakan ini dinyatakan sebagai
\begin{equation*}
E_{2} + \left( -2 \right) E_{1} = b_{2} + \left( -2 \right)b_{1}
\end{equation*}
atau
\begin{equation*}
2x-y + \left( -2 \right) \left(x+y \right) = 8 + \left( -2 \right) 7
\end{equation*}
Hal ini menghasilkan sistem kedua dalam Contoh \ref{example:effectsofelementaryoperation}, yaitu
\begin{equation*}
\begin{array}{c}
E_{1} = b_{1} \\
E_{2} + \left( -2 \right) E_{1} = b_{2} + \left( -2 \right) b_{1}
\end{array}
\end{equation*}

Dari titik ini, kita dapat
menemukan solusi sistem. Teorema \ref{thm:elementaryoperationsandsolns} memberi tahu kita bahwa solusi yang
kita temukan memang merupakan solusi sistem semula.

\ifdefined\showproofs
Sekarang kita akan membuktikan Teorema \ref{thm:elementaryoperationsandsolns}.

\begin{proof} 
\begin{enumerate}
\item
Pembuktian bahwa sistem \ref{system} dan \ref{thm1.9.1} memiliki
himpunan solusi yang sama adalah sebagai berikut. Misalkan $\left( x_{1},\cdots
,x_{n}\right) $ merupakan solusi $E_{1}=b_{1},E_{2}=b_{2}$. Kita ingin
menunjukkan bahwa pasangan ini merupakan solusi sistem pada \ref{thm1.9.1} di atas.
Hal ini jelas karena sistem pada \ref{thm1.9.1} merupakan sistem semula,
tetapi dicantumkan dalam urutan berbeda. Mengubah urutan tidak
memengaruhi himpunan solusi, sehingga $\left( x_{1},\cdots ,x_{n}\right) $ merupakan
solusi \ref{thm1.9.1}.

\item
Selanjutnya, kita ingin membuktikan bahwa sistem \ref{system} dan \ref{thm1.9.2} memiliki
himpunan solusi yang sama. Artinya, $E_{1}=b_{1},E_{2}=b_{2}$ memiliki
himpunan solusi yang sama dengan sistem $E_{1}=b_{1},kE_{2}=kb_{2}$ asalkan $k\neq
0 $. Misalkan $\left( x_{1},\cdots ,x_{n}\right) $ merupakan
solusi $E_{1}=b_{1},E_{2}=b_{2},$. Kita ingin menunjukkan bahwa pasangan ini merupakan solusi $E_{1}=b_{1},kE_{2}=kb_{2}$.
Perhatikan bahwa satu-satunya perbedaan di antara kedua sistem ini ialah bahwa sistem kedua
mengalikan persamaan $E_{2}=b_{2}$ dengan skalar $k$. Ingatlah bahwa ketika kedua ruas suatu
persamaan dikalikan dengan bilangan yang sama, kedua ruas tersebut tetap sama. Jadi, jika $\left( x_{1},\cdots ,x_{n}\right) $
merupakan solusi $E_{2}=b_{2}$, pasangan ini juga merupakan solusi $kE_{2}=kb_{2}$. Dengan demikian, $\left( x_{1},\cdots ,x_{n}\right) $ juga
merupakan solusi \ref{thm1.9.2}.

Demikian pula, misalkan $\left( x_{1},\cdots
,x_{n}\right) $ merupakan solusi $E_{1}=b_{1},kE_{2}=kb_{2}$. Selanjutnya kita dapat
mengalikan persamaan $kE_{2}=kb_{2}$ dengan skalar $1/k$, yang hanya mungkin karena kita mensyaratkan $k\neq
0$. Seperti di atas, tindakan ini mempertahankan kesamaan dan kita memperoleh persamaan $E_{2} = b_{2}$.
Jadi, $\left( x_{1},\cdots ,x_{n}\right)$ juga merupakan solusi $E_{1}=b_{1},E_{2}=b_{2}.$

\item
Terakhir, kita akan membuktikan bahwa sistem \ref{system} dan
\ref{thm1.9.3} memiliki himpunan solusi yang sama. Kita akan menunjukkan bahwa setiap
solusi $E_{1}=b_{1},E_{2}=b_{2}$ juga merupakan solusi
\ref{thm1.9.3}. Kemudian, kita akan menunjukkan bahwa setiap solusi \ref{thm1.9.3}
juga merupakan solusi $E_{1}=b_{1},E_{2}=b_{2}$. Misalkan $\left(
x_{1},\cdots ,x_{n}\right) $ merupakan solusi
$E_{1}=b_{1},E_{2}=b_{2}$. Secara khusus, pasangan ini memenuhi $E_{1} =
b_{1}$. Jadi, pasangan ini memenuhi persamaan pertama dalam \ref{thm1.9.3}.
Demikian pula, pasangan ini juga memenuhi $E_{2} = b_{2}$. Berdasarkan pembuktian
\ref{thm1.9.2}, pasangan ini juga memenuhi $kE_{1}=kb_{1}$. Perhatikan bahwa jika kita menjumlahkan
$E_{2}$ dan $kE_{1}$, hasilnya sama dengan $b_{2} + kb_{1}$. Oleh karena itu, jika
$\left( x_{1},\cdots ,x_{n}\right) $ memenuhi $E_{1}=b_{1},E_{2}=b_{2}$,
pasangan ini juga harus memenuhi $E_{2}+kE_{1}=b_{2}+kb_{1}$.

Sekarang misalkan $\left( x_{1},\cdots ,x_{n}\right)$ memenuhi sistem
$E_{1}=b_{1}, E_{2}+kE_{1}=b_{2}+kb_{1}$. Secara khusus, pasangan ini merupakan solusi $E_{1} = b_{1}$. Sekali lagi berdasarkan pembuktian \ref{thm1.9.2},
 pasangan ini juga merupakan solusi $kE_{1}=kb_{1}$. Jika kita mengurangkan besaran-besaran yang sama ini dari kedua ruas
$E_{2}+kE_{1}=b_{2}+kb_{1}$, kita memperoleh $E_{2}=b_{2}$, yang menunjukkan bahwa solusi tersebut juga memenuhi
$E_{1}=b_{1},E_{2}=b_{2}.$
\end{enumerate}
\end{proof}
\fi
Secara sederhana, teorema di atas menunjukkan bahwa operasi elementer tidak
mengubah himpunan solusi suatu sistem persamaan.

Sekarang kita akan melihat contoh suatu sistem dengan tiga persamaan dan tiga variabel.
Seperti pada contoh-contoh sebelumnya, tujuannya adalah menemukan nilai $x,y,z$ sedemikian sehingga setiap persamaan yang diberikan
terpenuhi ketika nilai-nilai tersebut disubstitusikan.

\begin{example}{ Menyelesaikan Sistem Persamaan dengan Operasi Elementer}\label{exa:solvingasystemwithelementaryops}
Temukan solusi sistem berikut,

\begin{equation}
\begin{array}{ccccccc}
x&+&3y&+&6z&=&25 \\
2x&+&7y&+&14z&=&58 \\
&&2y&+&5z&=&19
\end{array}
\label{solvingasystem1}
\end{equation}
\end{example}

\begin{solution}
Kita dapat menghubungkan sistem ini dengan Teorema \ref{thm:elementaryoperationsandsolns} di atas. Dalam hal ini, kita memiliki
\begin{equation*}
\begin{array}{c c}
E_{1} = x + 3y + 6z, & b_{1} = 25\\
E_{2} = 2x+7y+14z, & b_{2} = 58 \\
E_{3} = 2y+5z, & b_{3} = 19 
\end{array}
\end{equation*}
Teorema \ref{thm:elementaryoperationsandsolns} menyatakan bahwa jika kita melakukan operasi elementer pada sistem ini,
kita tidak akan mengubah himpunan solusinya. Oleh karena itu,
kita dapat menyelesaikan sistem ini menggunakan operasi elementer yang diberikan dalam Definisi \ref{def:elementaryoperations}.
Pertama, ganti persamaan kedua dengan hasil penambahan $\left( -2\right)$
kali persamaan pertama ke persamaan kedua. Hal ini menghasilkan sistem
\begin{equation}
\begin{array}{c}
x+3y+6z=25 \\
y+2z=8 \\
2y+5z=19
\end{array}
\label{solvingasystem2}
\end{equation}
Sekarang, ganti persamaan ketiga dengan hasil penambahan $\left( -2\right) $ kali
persamaan kedua ke persamaan ketiga. Hal ini menghasilkan sistem
\begin{equation}
\begin{array}{c}
x+3y+6z=25 \\
y+2z=8 \\
z=3
\end{array}
\label{solvingasystem3}
\end{equation}
Pada tahap ini, kita dapat menemukan solusinya dengan mudah. Ambil $z=3$ dan substitusikan kembali nilai ini ke dalam
 persamaan sebelumnya untuk mencari $y$, lalu lakukan hal serupa untuk mencari $x$.
\begin{equation*}
\begin{array}{c}
x + 3y + 6 \left(3 \right) = x + 3y + 18 = 25\\
y + 2 \left(3 \right) = y + 6 = 8 \\
z = 3
\end{array}
\end{equation*}
Persamaan kedua kini menjadi
\begin{equation*}
y+6=8
\end{equation*}
Dari persamaan ini Anda dapat melihat bahwa $y = 2$. Oleh karena itu, kita dapat menyubstitusikan nilai tersebut ke dalam persamaan pertama sebagai berikut:
\begin{equation*}
x + 3 \left(2 \right) + 18 = 25
\end{equation*}
Dengan menyederhanakan persamaan ini, kita menemukan bahwa $x=1$.
Jadi, solusi sistem ini adalah $\left( x,y,z \right) = \left( 1,2,3 \right)$.
Proses ini disebut \textbf{substitusi balik}. \index{substitusi balik}

Sebagai alternatif, pada \ref{solvingasystem3} Anda dapat melanjutkan sebagai berikut. Tambahkan $\left(
-2\right) $ kali persamaan ketiga ke persamaan kedua, lalu tambahkan $\left(
-6\right) $ kali persamaan kedua ke persamaan pertama. Hal ini menghasilkan
\begin{equation*}
\allowbreak
\begin{array}{c}
x+3y=7 \\
y=2 \\
z=3
\end{array}
\end{equation*}
Sekarang tambahkan $\left( -3\right) $ kali persamaan kedua ke persamaan pertama. Hal ini menghasilkan
\begin{equation*}
\allowbreak
\begin{array}{c}
x=1 \\
y=2 \\
z=3
\end{array}
\end{equation*}
suatu sistem yang memiliki himpunan solusi yang sama dengan sistem semula. Cara ini
menghindari substitusi balik dan menghasilkan himpunan solusi yang sama. Anda dapat memilih metode yang lebih disukai karena keduanya menghasilkan solusi yang benar,
$\left( x,y,z \right) = \left(1,2,3\right)$.
\end{solution}
